\documentclass[10pt,a4paper]{amsart}
\usepackage[margin=19mm]{geometry}
\usepackage{amsmath,amssymb,mathtools}
\usepackage[scr=rsfs,cal=euler]{mathalfa}
\usepackage{xfrac}
\usepackage[unicode,hidelinks,bookmarks=false]{hyperref}
\newcommand{\ie}{i.e.\ }
\newcommand{\cf}{cf.\ }
\newcommand{\resp}{respectively\ }
\newcommand{\Subsection}{\subsection}
\providecommand{\nobreakdash}{\nobreak\hbox{-}}
\setlength{\emergencystretch}{2em}
\multlinegap0pt
\usepackage{bm}
\usepackage{bbm}
\usepackage{dsfont}
\usepackage{xspace}
\usepackage{cancel}
\usepackage{mathtools}
\usepackage{amsthm}
\makeatletter
\@ifundefined{theorem}{\newtheorem{theorem}{Theorem}}{}
\@ifundefined{lemma}{\newtheorem{lemma}[theorem]{Lemma}}{}
\makeatother
\newcommand{\Lipz}{\operatorname{Lip}_0}
\newcommand{\N}{\mathbb{N}}
\newcommand{\Z}{\mathbb{Z}}
\newcommand{\absFunc}{\abs{\: \cdot \: }}
\newcommand{\abs}[1]{\left| #1 \right|}
\newcommand{\eps}{\varepsilon}
\newcommand{\norm}[1]{\left\| #1 \right\|}
\newcommand{\ord}{\operatorname{ord}}
\title[Almost-invariant elements of group actions on Lip$\_0$ spaces]{Almost-invariant elements of group actions on Lip$\_0$ spaces}
\author{Tom\'a\v{s} Raunig}
\date{Source: arXiv:2506.06770v1 (CC BY 4.0)}
\begin{document}
\maketitle

\noindent\emph{Source and licence.} Almost-invariant elements of group actions on Lip$_0$ spaces. Version arXiv:2506.06770v1 (CC BY 4.0), distributed under the Creative Commons Attribution 4.0 International licence, as recorded in the arXiv licence metadata for this exact version and shown on the abstract page https://arxiv.org/abs/2506.06770v1. Full rights notice: licenses/arxiv-2506.06770-CC-BY-4.0.txt.\par\smallskip

\noindent\textit{Editorial scope.} Counted unit: the Lemma whose source label is \texttt{lma:badZ} in the article (source0.tex lines 701--703, source label \texttt{lma:badZ}), delivered together with the article's own complete original proof of it (source0.tex lines 704--764, one \texttt{proof} environment of 61 lines). No mathematics is written, repaired or re-derived: statement and proof are the article's own text line for line. Required context retained uncounted: none. Every definition, display and label of those lines is retained unchanged. Omitted from this package and not redistributed: 1--700, 765--950. Nothing inside a retained excerpt is omitted or altered: every retained body is delivered exactly as the article's lines are written, so no omission receipt of any kind accompanies this package. Citations: the retained excerpts carry no citation command at all, so no bibliography entry of the article is transcribed and no reference is redistributed. Reference adaptation: none was needed and none was made; every reference inside the delivered unit resolves inside it, so no unresolved reference or citation is produced. Printed slips kept literal: no printed word, symbol, hypothesis or number of the counted statement or of its proof was altered. Layout and class: the article's own class is \texttt{amsart} with packages including inputenc, babel, a4wide, amsaddr, amsmath, amssymb, amsthm, cleveref, xcolor, enumitem, todonotes, url; the wrapper is the lane's pinned \texttt{amsart} editorial wrapper at 10pt on A4, which restates the theorem-like environments the retained text needs and adds the article's own macro definitions that the retained text needs (\texttt{\textbackslash Lipz}, \texttt{\textbackslash N}, \texttt{\textbackslash Z}, \texttt{\textbackslash abs}, \texttt{\textbackslash absFunc}, \texttt{\textbackslash eps}, \texttt{\textbackslash norm}, \texttt{\textbackslash ord}), with extra wrapper packages (bm, bbm, dsfont, xspace, cancel, mathtools). The delivered numbering is the unit's own. Assets: no retained span contains a figure environment, a graphics-inclusion command, a PDF-page inclusion, a table or a TikZ picture; no image file is copied into this package, the package declares no asset dependency and no image file is redistributed with it. Licence: version arXiv:2506.06770v1 of this article is distributed under the Creative Commons Attribution 4.0 International licence, as recorded in the arXiv licence metadata for this exact version and shown on the abstract page https://arxiv.org/abs/2506.06770v1. A full rights notice is delivered at licenses/arxiv-2506.06770-CC-BY-4.0.txt. Characters: every retained excerpt is pure ASCII, so the delivered engine, XeLaTeX with its default Latin Modern text font, reports no missing character.
	\begin{lemma} \label{lma:badZ}
		Let $G$ be a~cyclic group equipped with a~left-invariant pseudometric $d$ and the action by translations on itself. If $(G, d)$ is not bi-Lipschitz equivalent to $(\Z, \abs{\cdot})$, then for every $\delta > 0$ and $\delta$-invariant $f \in \Lipz(G)$ it is the case that $\norm{f} \leq \delta$.
	\end{lemma}

	\begin{proof}
		For the cyclic (and hence commutative) group $G$ we adopt the additive notation. For the ease of notation, we assume $G$ to be either $\Z$ or $\Z_n$ for some $n \in \N$. If we have the latter, all operations are to be understood as being taken modulo $n$.
		
		First, since $G$ acts by isometries, we obtain
		\begin{equation*}
			\forall k, l \in \Z, k \geq l \colon d(k, l) = d(k-l, 0) \leq d(1, 0) + d(1, 2) + \cdots + d(k-l, k-l-1) = \abs{k-l}d(1, 0),
		\end{equation*}
		i.e. $d \leq d(0,1)\absFunc$. Hence we can without loss of generality assume that $d \leq \absFunc$ by considering the function $f/d(0,1)$ and metric $d/d(0,1)$. We assume that $d(0, 1) \neq 0$, because if $d(0, 1) = 0$, then the distance of any two points is zero and hence the only Lipschitz function on $G$ is constant zero.
		
		For a~contradiction, assume that $\norm{f}_{(\Z, d)} > \delta$. Then there are $\eps > 0$ and $l, l' \in G, l > l'$ such that $\abs{f(l) - f(l')}/d(l, l') > \delta + \eps$. Potentially passing to $-f$, we may, without loss of generality, assume that $(f(l) - f(l'))/d(l, l') > \delta + \eps$. Denote $k = l - l'$ and let $n \in \N$. Using the assumption that $\norm{gf - f} \leq \delta$ for $g = l' - nk$ we obtain
		\begin{equation*}
			\abs{(f((n+1)k) - f(nk)) - (f(l) - f(l'))}
			= \abs{(gf(l) - f(l)) - (gf(l') - f(l'))}
			\leq \delta d(l, l')
		\end{equation*}
		and hence
		\begin{align*}
			f((n+1)k) - f(nk)
			&= f(l) - f(l') - (f((n+1)k) - f(nk)) - (f(l) - f(l'))\\
			&> (\delta + \eps) d(l, l') - \delta d(l, l')
			= \eps d(l, l')
			= \eps d(k, 0).
		\end{align*}
		That is, the sequence $(f(nk))_{n=1}^\infty$ is monotone and with each step increases by at least $\eps d(k, 0)$. Thus, if we put $\alpha = \eps d(k, 0)$, then, because $f$ is Lipschitz and $f(l) \neq f(l')$, $\alpha > 0$ and 
		\begin{equation} \label{eqn:badZ}
			f(nk) > \alpha n, n \in \N.
		\end{equation}
		
		Now we discern three cases. First we deal with the easiest case when $d$ is truly a~pseudometric and not a~metric. Then, using left invariance, we have some $j \in G$ such that $d(0, j) = 0$ and hence also $d(0, jn) = 0$ for all $n \in \N$. But this is in contradiction with~\eqref{eqn:badZ}, because
		\begin{equation*}
			0 < \alpha j < f(jk) = f(jk) - f(0) \leq \norm{f} d(jk, 0) = 0.
		\end{equation*}
		
		Assume $d$ is a~metric, but $G$ is not (algebraically) isomorphic to $\Z$. Under this assumption, $\ord_G(k) < \infty$ and so, in \eqref{eqn:badZ} we may take $n = \ord_G(k)$, yielding
		\begin{equation*}
			0 = f(0) = f(\ord_G(k) k) > \alpha \ord_G(k),
		\end{equation*}
		which is, of course, a~contradiction, because, by definition of $k$, $\ord_G(k) \neq 0$.
		
		When $G$ is isomorphic to $\Z$, one must make a~slightly finer argument, starting with the following estimate
		\begin{equation} \label{eqn:KYD_ZonZ}
			\forall n \in \N \colon L(f) \geq \frac{f(nk) - f(0)}{d(nk, 0)}
			> \frac{\alpha n}{d(nk, 0)}
			= \frac{\alpha}{k} \cdot \frac{nk}{d(nk, 0)}.
		\end{equation}
		Choose $m \in \N, m > kL(f) /\alpha$. Using triangle inequality we obtain
		\begin{equation*}
			\frac{k}{\alpha} L(f)
			< m
			< \frac{k_m}{d(k_m, 0)}
			= \frac{kk_m}{kd(k_m, 0)}
			\leq \frac{kk_m}{d(kk_m, 0)}.
		\end{equation*}
		We finish by setting $n = k_m$ in~\eqref{eqn:KYD_ZonZ}:
		\begin{equation*}
			L(f) >
			\frac{\alpha}{k} \cdot \frac{k_m k}{d(k_m k, 0)}
			> \frac{\alpha}{k} \cdot \frac{k}{\alpha} L(f),
		\end{equation*}
		which is a~contradiction.
	\end{proof}

\end{document}
